% !TeX program = pdflatex
% !TeX TXS-program:compile = txs:///pdflatex/[-shell-escape]
% !TeX TXS-program:pdflatex = pdflatex -synctex=1 -interaction=nonstopmode --shell-escape %.tex

\section{Template}
\label{chap:template}
\setCustomHeader{{\color{csecPurple}\faBrain}\ \textbf{Playbook}}{Knowledge Base \& Bitácora}{\today}

Este capítulo es un ejemplo de como se tienen que realizar las demás secciones del libro
% ================================================================================
% 1. Metodología de Registro y Fuentes de Estudio
% ================================================================================
\subsection{Metodología de Registro y Fuentes de Consulta}
Para que el conocimiento perdure y sea consultable de manera inmediata en auditorías reales o exámenes prácticos (como OSCP o CPTS), cada concepto registrado debe responder a tres preguntas fundamentales:
\begin{enumerate}
	\item \textbf{¿Qué aprendí?} (Concepto técnico, mecanismo de funcionamiento, payload o comando clave).
	\item \textbf{¿Dónde lo aprendí?} (Libro con autor y capítulo, sitio web/documentación con URL, o máquina vulnerable con enlace interno).
	\item \textbf{¿Cómo se aplica?} (Comando reproducible, script o prueba de concepto).
\end{enumerate}

\begin{infobox}[Badges de Fuentes Disponibles]
	Para estandarizar el origen del conocimiento, dispones de la macro \texttt{\textbackslash sourceBadge\{\ldots\}} y el comando \texttt{\textbackslash resourceItem\{Tipo\}\{Título/Autor\}\{Detalle/Enlace\}}:
	\begin{center}
		\small
		\sourceBadge{Libro} \hspace{4pt}
		\sourceBadge{Web} \hspace{4pt}
		\sourceBadge{Docs} \hspace{4pt}
		\sourceBadge{Curso} \hspace{4pt}
		\sourceBadge{Lab} \hspace{4pt}
		\sourceBadge{Writeup} \hspace{4pt}
		\sourceBadge{Video} \hspace{4pt}
		\sourceBadge{Paper}
	\end{center}
\end{infobox}


\subsection{Ejemplo de algunas cajas}

\begin{learningbox}[Conceptos Clave Aprendidos: Bash Avanzado]
	\begin{itemize}[leftmargin=*]
		\item \textbf{Redirección y Descriptores de Ficheros:} Control preciso de \texttt{stdin} (0), \texttt{stdout} (1) y \texttt{stderr} (2). Uso de \texttt{2>\&1} para unificar flujos y \texttt{/dev/null} para silenciar ruido en bucles.
		\item \textbf{One-Liners de Reconocimiento:} Uso de combinaciones de \texttt{curl}, \texttt{grep}, \texttt{sed}, \texttt{awk} y \texttt{xargs} para extraer subdominios o endpoints activos en segundos sin herramientas pesadas.
		\item \textbf{Subshells y Captura de Salida:} Diferencia entre asignación síncrona mediante \texttt{\$(comando)} y ejecución en subshell background con \texttt{\&} junto al comando \texttt{wait}.
		\item \textbf{Manejo de Señales (\texttt{trap}):} Limpieza automática de ficheros temporales y sockets al pulsar \texttt{Ctrl+C} (\texttt{SIGINT}) en scripts interactivos.
	\end{itemize}
\end{learningbox}

\begin{sourcebox}[Fuentes y Bibliografía Utilizadas]
	\begin{itemize}[leftmargin=*]
		\resourceItem{Libro}{The Linux Command Line: A Complete Introduction (William Shotts)}{Capítulos 20 al 35: Entornos shell, variables, control de flujo y scripts robustos}
		\resourceItem{Libro}{Wicked Cool Shell Scripts (Dave Taylor \& Brandon Perry)}{2ª Edición — Scripts prácticos para administración de sistemas y seguridad}
		\resourceItem{Web}{Bash Hackers Wiki — Guía profunda de sintaxis y buenas prácticas}{\url{https://wiki.bash-hackers.org}}
		\resourceItem{Web}{ExplainShell — Análisis visual e interactivo de comandos y flags}{\url{https://explainshell.com}}
		\resourceItem{Docs}{GNU Bash Reference Manual Oficial}{\url{https://www.gnu.org/software/bash/manual/}}
	\end{itemize}
\end{sourcebox}


\begin{bashbox}{Snippet Útil: Escaneo Rápido de Hosts Vía Ping en Bash}
	#!/usr/bin/env bash
	# Sweep rápido de red /24 en paralelo usando subshells
	SUBNET="10.10.10"
	echo "[*] Escaneando subred ${SUBNET}.0/24..."

		for ip in $(seq 1 254); do
	(ping -c 1 -W 1 "${SUBNET}.${ip}" &>/dev/null && echo "[+] Host activo: ${SUBNET}.${ip}") &
	done
	wait
	echo "[*] Barrido completado."
\end{bashbox}
